% Part: model-theory
% Chapter: models-of-arithmetic
% Section: standard-models

\documentclass[../../../include/open-logic-section]{subfiles}

\begin{document}

\olfileid{mod}{mar}{stm}
\olsection{Model-Model Standar Aritmetika}

Basa aritmetika~$\Lang{L_A}$
mesthi dimaksudake kanggo
ngrembug wilangan, mliginé
wilangan asli. Mula model
standar~$\Struct{N}$ iku
istimewa: iki model kang
arep dirembug. Nanging
ing logika, kang kerep
dadi kawigaten kita yaiku
sipat struktural, lan
sembarang rong !!{structure}s
kang isomorfis nduweni
sipat-sipat iku padha.
Mula kita bisa luwih luwes
lan nganggep sembarang
!!{structure} kang isomorfis
karo~$\Struct{N}$
minangka standar.

\begin{defn}
!!{structure} kanggo
$\Lang{L_A}$ iku
\emph{standar} yen isomorfis
karo~$\Struct{N}$.
\end{defn}

\begin{prop}
\ollabel{prop:standard-domain}
Yen !!a{structure}~$\Struct{M}$
standar, ranahé iku
himpunan nilai angka-angka
standar, yaiku
\[
\Domain{M} = \Setabs{\Value{\num{n}}{M}}{n \in \Nat}
\]
\end{prop}

\begin{proof}
Cetha yen saben
$\Value{\num{n}}{M} \in \Domain{M}$.
Kita mung kudu nuduhake
yen saben $x \in \Domain{M}$
padha karo $\Value{\num{n}}{M}$
kanggo sawenehing~$n$.
Amarga $\Struct{M}$ standar,
struktur iku isomorfis
karo~$\Struct{N}$.
Upamakna $g\colon
\Nat \to \Domain{M}$
iku isomorfisme.
Mula $g(n) =
g(\Value{\num{n}}{N}) =
\Value{\num{n}}{M}$.
Nanging kanggo saben
$x \in \Domain{M}$,
ana~$n
\in \Nat$ saéngga
$g(n) = x$,
amarga $g$ !!{surjective}.
\end{proof}

\begin{explain}
Yen struktur~$\Struct{M}$
kanggo $\Lang{L_A}$ standar,
kabeh unsur ing
!!{domain}-é bisa dijenengi
nganggo angka standar
$\num{0}$, $\num{1}$,
$\num{2}$, \dots,
yaiku term $\Obj{0}$,
$\Obj{0}'$, $\Obj{0}''$,
lan sateruse. Mesthi
iki ora ateges
!!{element}s saka
$\Domain{M}$ \emph{pancen}
wilangan, nanging mung
ateges kita bisa
nunjuk unsur-unsur mau
kanthi cara padha
kaya nuduhake wilangan
ing~$\Domain{N}$.
\end{explain}

\begin{prob}
Tuduhna yen kosok baline
\olref[mod][mar][stm]{prop:standard-domain}
ora bener, yaiku wenehana
tuladha !!a{structure}~$\Struct{M}$
kanthi $\Domain{M} =
\Setabs{\Value{\num{n}}{M}}{n \in \Nat}$
kang ora isomorfis
karo~$\Struct{N}$.
\end{prob}

\begin{prop}
\ollabel{prop:thq-standard}
Yen $\Sat{M}{\Th{Q}}$,
lan $\Domain{M} =
\Setabs{\Value{\num{n}}{M}}{n
  \in \Nat}$,
mula $\Struct{M}$ standar.
\end{prop}

\begin{proof}
Kita kudu nuduhake yen
$\Struct{M}$ isomorfis
karo~$\Struct{N}$.
Gatekna fungsi
$g\colon \Nat \to \Domain{M}$
kang ditetepake dening
$g(n) = \Value{\num{n}}{M}$.
Miturut hipotesis,
$g$ !!{surjective}.
Fungsi iku uga
!!{injective}:
$\Th{Q} \Proves
\eq/[\num{n}][\num{m}]$
manawa $n \neq m$.
Mula, amarga
$\Sat{M}{\Th{Q}}$,
$\Sat{M}{\eq/[\num{n}][\num{m}]}$
manawa $n \neq
m$. Dadi, yen
$n \neq m$,
mula $\Value{\num{n}}{M}
\neq \Value{\num{m}}{M}$,
yaiku $g(n) \neq g(m)$.

Kita uga kudu mriksa yen
$g$ iku isomorfisme.
\begin{enumerate}
\item Kita nduweni
  $g(\Assign{\Obj{0}}{N}) = g(0)$
  amarga
  $\Assign{\Obj{0}}{N} = 0$.
  Miturut definisi~$g$,
  $g(0) = \Value{\num{0}}{M}$.
  Nanging $\num{0}$
  iku mung $\Obj{0}$,
  lan nilai term kang
  arupa !!a{constant}
  ditemtokake dening apa
  kang diwenehake
  !!{structure} marang
  !!{constant} iku,
  yaiku
  $\Value{\Obj{0}}{M} =
  \Assign{\Obj{0}}{M}$.
  Mula kita nduweni
  $g(\Assign{\Obj{0}}{N}) =
  \Assign{\Obj{0}}{M}$
  kaya kang dibutuhake.
\item $g(\Assign{\prime}{N}(n)) =
  g(n+1)$, amarga
  $\prime$ ing
  $\Struct{N}$ iku
  fungsi penerus ing~$\Nat$.
  Banjur $g(n+1) =
  \Value{\num{n+1}}{M}$
  miturut definisi~$g$.
  Nanging $\num{n+1}$
  iku term kang padha karo
  $\num{n}'$, mula
  $\Value{\num{n+1}}{M} =
  \Value{\num{n}'}{M}$.
  Miturut definisi
  fungsi nilai, iki
  $= \Assign{\prime}{M}(
  \Value{\num{n}}{M})$.
  Amarga
  $\Value{\num{n}}{M} = g(n)$,
  kita entuk
  $g(\Assign{\prime}{N}(n)) =
  \Assign{\prime}{M}(g(n))$.
\item $g(\Assign{+}{N}(n,m)) =
  g(n+m)$, amarga
  $+$ ing $\Struct{N}$
  iku fungsi panjumlahan
  ing~$\Nat$.
  Banjur $g(n+m) =
  \Value{\num{n+m}}{M}$
  miturut definisi~$g$.
  Nanging
  $\Th{Q} \Proves
  \num{n+m} =
  (\num{n} + \num{m})$,
  mula $\Value{\num{n+m}}{M} =
  \Value{\num{n}+\num{m}}{M}$.
  Miturut definisi
  fungsi nilai, iki
  $=
  \Assign{+}{M}(
  \Value{\num{n}}{M},
  \Value{\num{m}}{M})$.
  Amarga
  $\Value{\num{n}}{M} = g(n)$
  lan $\Value{\num{m}}{M} = g(m)$,
  kita entuk
  $g(\Assign{+}{N}(n, m)) =
  \Assign{+}{M}(g(n), g(m))$.
\item $g(\Assign{\times}{N}(n, m)) =
  \Assign{\times}{M}(g(n), g(m))$:
  Gladhen.
\item $\tuple{n,m} \in
  \Assign{<}{N}$ yen lan mung yen
  $n < m$. Yen $n < m$,
  mula $\Th{Q} \Proves
  \num{n} < \num{m}$,
  lan uga $\Sat{M}{\num{n} <
  \num{m}}$.
  Mula $\tuple{
  \Value{\num{n}}{M},
  \Value{\num{m}}{M}} \in
  \Assign{<}{M}$,
  yaiku $\tuple{g(n), g(m)}
  \in \Assign{<}{M}$.
  Yen $n
  \not< m$,
  mula $\Th{Q} \Proves
  \lnot \num{n} < \num{m}$,
  lan akibate
  $\Sat/{M}{\num{n} < \num{m}}$.
  Mula, kaya ing ndhuwur,
  $\tuple{g(n), g(m)} \notin
  \Assign{<}{M}$.
  Dadi kita entuk:
  $\tuple{n,m} \in
  \Assign{<}{N}$ yen lan mung yen
  $\tuple{g(n), g(m)} \in
  \Assign{<}{M}$.
\end{enumerate}
\end{proof}

\begin{explain}
Fungsi~$g$ iku cara
paling cetha kanggo
netepake pemetaan saka
$\Nat$ menyang ranah
sembarang !!{structure}~$\Struct{M}$
liyane kanggo
$\Lang{L_A}$, amarga
saben $\Struct{M}$
kaya mangkono ngemot
!!{element}s kang dijenengi
dening $\num{0}$,
$\num{1}$, $\num{2}$,
lan sateruse.
Mula ora nggumunake yen
$\Struct{M}$ ndadekake
paling ora sawenehing
pratelan dhasar ngenani
$\num{n}$ bener
kanthi cara padha
kaya~$\Struct{N}$,
lan $g$ uga bijektif,
mula $g$ dadi isomorfisme.
Nyatane, yen
$\Domain{M}$ ora ngemot
!!{element}s liyane
ngluwihi kang dijenengi
dening $\num{n}$,
iku siji-sijine.
\end{explain}

\begin{prop}
\ollabel{prop:thq-unique-iso}
Yen $\Struct{M}$ standar,
mula $g$ saka bukti
\olref{prop:thq-standard}
iku siji-sijine isomorfisme
saka $\Struct{N}$
menyang~$\Struct{M}$.
\end{prop}

\begin{proof}
Upamakna $h\colon
\Nat \to \Domain{M}$
iku isomorfisme
antarane $\Struct{N}$
lan~$\Struct{M}$.
Kita nuduhake $g = h$
kanthi induksi tumrap~$n$.
Yen $n = 0$, mula
$g(0) = \Assign{\Obj{0}}{M}$
miturut definisi~$g$.
Nanging amarga $h$
isomorfisme,
$h(0) =
h(\Assign{\Obj{0}}{N}) =
\Assign{\Obj{0}}{M}$,
mula $g(0) = h(0)$.

Saiki gatekna kasus
$n+1$. Kita nduweni
\begin{align*}
  g(n+1) & = \Value{\num{n+1}}{M} \text{ miturut definisi~$g$}\\
  & = \Value{\num{n}'}{M} \text{ amarga $\num{n+1}\ident \num{n}'$}\\
  & = \Assign{\prime}{M}(\Value{\num{n}}{M})
    \text{ miturut definisi $\Value{t'}{M}$}\\
  & = \Assign{\prime}{M}(g(n))  \text{ miturut definisi~$g$}\\
  & = \Assign{\prime}{M}(h(n)) \text{ miturut hipotesis induksi}\\
  & = h(\Assign{\prime}{N}(n)) \text{ amarga $h$ iku isomorfisme}\\
  & = h(n+1)
\end{align*}
\end{proof}

\begin{explain}
Kanggo sembarang himpunan
!!{denumerable}~$M$,
ana !!a{bijection}
antarane $\Nat$ lan
$M$, mula saben himpunan~$M$
kaya mangkono bisa dadi
!!{domain} model
standar~$\Struct{M}$.
Nyatane, sawisé milih
objek $z \in M$
lan fungsi $s$
kang cocog minangka
$\Assign{\Obj{0}}{M}$
lan $\Assign{\prime}{M}$,
interpretasi $+$,
$\times$, lan $<$
wis katetepake.
Mung fungsi~$s\colon
M \to M \setminus \{z\}$
kang !!{injective}
lan !!{surjective}
kang bisa dadi
$\Assign{\prime}{M}$
ing model standar.
Jajaran asil $s$
ora bisa ngemot~$z$,
amarga yen ngemot,
$\lforall[x][
\eq/[\Obj 0][x']]$
bakal kleru.
!!{sentence} iku
bener ing~$\Struct{N}$,
mula $\Struct{M}$
uga kudu ndadekake
ukara iku bener.
Fungsi~$s$ kudu
!!{injective},
amarga fungsi penerus
$\Assign{\prime}{N}$
ing~$\Struct{N}$
uga mangkono, lan
sipat !!{injective}
saka $\Assign{\prime}{N}$
bisa dinyatakake
dening !!a{sentence}
kang bener ing~$\Struct{N}$.
Fungsi iku kudu
!!{surjective}
amarga yen ora,
ana sawenehing
$x \in M \setminus \{z\}$
kang ora ana ing
jajaran asil saka~$s$;
kanthi mangkono
!!{sentence}
$\lforall[x][(
\eq[x][\Obj 0] \lor
\lexists[y][\eq[y'][x]])]$
bakal kleru ing~$\Struct{M}$,
dene ukara iku bener
ing~$\Struct{N}$.
Cathetan koreksi sumber OLPL-096:
surjektivitas gegayutan
karo jajaran asil fungsi,
dudu ranah inpute.
\end{explain}


\end{document}
